\documentclass[12pt,a4paper]{article}
\usepackage[UTF8]{ctex}
\usepackage{amsfonts,amsmath,amsthm,amssymb}
\usepackage{sfmath}
\usepackage{hyperref}
\renewcommand{\familydefault}{\sfdefault}
\usepackage{geometry}
\geometry{left=0cm,right=0cm,top=0cm,bottom=0.1cm, footskip=0cm, includefoot=true}
%\usepackage[contents = \thepage, color=yellow, angle=0, scale = 20]{background}
\pagestyle{empty}
%\usepackage{fancyhdr}
%\pagestyle{fancy}
%\fancyhead{}
%\fancyfoot{}
%\cfoot{}
%\lfoot{{\small \color{red} \slshape \thepage 页}}
%\fancyfoot[R]{\thepage}

\usepackage[all]{xy}
\usepackage{color}

\renewenvironment{proof}{{\noindent\it\textbf{Proof}}\quad\it}{\hfill$\square$}

%\theoremstyle{definition}

\newtheorem{thm}[equation]{Theorem}
\newtheorem{cor}[equation]{Corollary}
\newtheorem{lem}[equation]{Lemma}
\newtheorem{prop}[equation]{Proposition}
\newtheorem{defn}[equation]{Definition}
\newtheorem{rem}[equation]{Remark}

\title{last Kervaire}
\author{}

\setcounter{secnumdepth}{6}

\begin{document}

\sffamily
\maketitle
\thispagestyle{empty}

\abstract{}
\tableofcontents

\large

\begin{sloppy}
\clubpenalty=-100
\widowpenalty=-100



\section{Adams谱序列}

\subsection{Adams 谱序列构造}

\paragraph{公理} 对任意谱形$X$, $Y$, 存在Adams 谱序列 $Adams(X,Y)$

\paragraph{公理}Adams谱序列是$\mathcal{S}^{op}\times \mathcal{S}$到双分次的谱序列范畴的函子

\paragraph{谱形的有限性}

\paragraph{定义} $\mathcal{S}^{fin}$ 为包含球谱的最小的全子范畴，在suspension、desuspension、cofiber操作下封闭

\paragraph{定理}  $Z\in\mathcal{S}^{fin}$ 当且仅当$X$可以通过有限步如下的操作得到：\begin{enumerate}
    \item suspension/desuspension
    \item taking cofiber under a map from the sphere spectrum
\end{enumerate}

\subsection{谱形的连通性}
\paragraph{定义}$X$是$n$-连通的，若对所有$i\leq n$,$\pi_i(X)=0$

\paragraph{定理sorry}
球谱是 $(-1)$-连通的

\paragraph{定义}X is bounded below iff there is some $n\in \mathbb{Z}$, such that X is n-connected

\paragraph{定理sorry} 若X是n-连通，则当$i\leq n$时，$H_*(X;\mathbb{F}_2) = 0$

\subsection{finite type}
\paragraph{定义} a spectrum $X$ is finite type iff
for any $n$, there is some spectrum $X_n\in \mathcal{S}^{fin}$, such that there is a map
$X_n\xrightarrow{}X$, whose cofiber is $n$-connected

\paragraph{定理sorry}$X$ finite type当且仅当
\begin{enumerate}
    \item $X$ is bounded below
    \item for any $i$, $\pi_i(X)$ is finitely generated as abelian group
\end{enumerate}

\paragraph{定理sorry}
$X$ finite type， 则$H_i（X;\mathbb{F}_2）$ 均为有限维

\subsection{Adams 谱序列的收敛性}
本subsection均假设$X\in \mathcal{S}^{fin}$, $Y\in \mathcal{S}$, $Y$ finite type。
\paragraph{Adams filtration的定义}

$f\in [X,Y]$ has Adams filtration $\geq n$ 当且仅当
存在$X_1,\dots, X_{n-1} \in \mathcal{S}^{fin}$, 使得$f$等于如下的复合
$$X\xrightarrow[]{f_1}X_1\xrightarrow{f_2}X_2\xrightarrow[]{f_3}\dots \xrightarrow[]{f_{n-1}}X_{n-1}\xrightarrow[]{f_n}Y$$
且$f_i$诱导的mod 2同调的映射为0

\paragraph{定义}称$f\in [X,Y]$的Adams filtration为无穷，若对所有n，$f$的Adams filtration $\geq n$

\paragraph{定理}
若$f$是奇数阶，则f的Adams filtration为无穷

\paragraph{定理sorry}
若$f$不是奇数阶，则f的Adams filtration不是无穷

\paragraph{定理sorry}
Adams谱序列$Adams(X,Y)$弱收敛到$[X,Y]$, 即Adams filtration的graded piece与Adams 谱序列的$E_\infty$项中相应的群之间存在自然同构

\paragraph{定理}
$[X,Y]/\text{所有奇数阶元}$上面具有诱导的Adams filtration，并且这时filtration is separated

\paragraph{Adams谱序列的boundedness}

若$X$ finite type

\paragraph{定理 sorry}
假设$H_i(X;\mathbb{Q}) = 0$, 则$\pi_i(X)/\text{所有奇数阶元}$上面的Adams filtation是bounded

\paragraph{定理sorry}
$\pi_i(X)/\text{所有奇数阶元}$上面的Adams filtation
满足所有的subquotient都为有限群


\begin{thebibliography}{99}


\end{thebibliography}

\end{sloppy}
\end{document}
